% =============================================================================
% DRAFT — Section 4 v3 + new Appendix items
% Replaces lines 563-818 in main.tex (the current Section 4 case study)
% AND adds appendix items at the end (Eval-v2 details, IC-trap, Agent traces)
% =============================================================================

\section{Case Study: Negative Knowledge in a Multi-Agent Research Loop}
\label{sec:case-study}

We test structured negative knowledge in a real research workflow on the
coupled Burgers-swept-KdV (BKdV) system of Holm et~al.\ (2025). Unlike
benchmark evaluation, this setting has no closed-form reference; ``success''
is whether the team produces a defensible physical observation. Failures
take three forms: numerical blow-up, runs that finish but miss the phenomenon,
and runs that are numerically valid but scientifically uninformative. The
case study tests whether bounded negative knowledge---with explicit depth
structure---helps a multi-agent team navigate this richer failure space.

The BKdV system couples a Burgers-like current $u$ to a KdV-like wave $v$,
\begin{align*}
  u_t + 3 u u_x &= -\partial_x\!\left(3 v^2 + \gamma v_{xx}\right), \\
  v_t + 6 v v_x + \gamma v_{xxx} &= -\partial_x(u v).
\end{align*}
The reduction $m := u - v^2/2 = 0$ collapses to a Gardner equation; the three
limits (Burgers, KdV, Gardner) are known while the coupled regime is partially
open. The case study runs in two stages: \textbf{Stage~1} produces a structured
knowledge bank from \emph{multi-round} stress tests on the coupled system
itself; \textbf{Stage~2} uses that bank to study three coupled-system
sub-questions under four memory conditions and a strict
progressive-complexity discipline. Full task specifications, the
bank inventory, the eval-version derivation, and the agent trace catalogue
are in
\cref{appendix:bkdv-setup,appendix:bkdv-bank,appendix:eval-v2,appendix:agent-traces}.

\subsection{Stage~1: Multi-round stress tests with depth-tagged curation}
\label{sec:pde-stage1}

\paragraph{From single-shot logs to depth-tagged knowledge.} Earlier rounds of
this case study produced a 30-entry bank from \emph{single-shot} stress tests
on isolated component PDEs. We found this insufficient: a one-round failure
is at most an error message---an agent at Stage~2 reproduces equivalent
information after one round on their own. Knowledge that genuinely helps must
capture what \emph{persists} across multiple attempts: a path-level closure
that the agent at Stage~2 would not reproduce in their own three-round
budget. We therefore reorganise Stage~1 into seven \emph{multi-round}
research-program stress tests (BKdV-S1 through BKdV-S7) on the coupled
system, each conducted under the Research Graph protocol with three Experiment
rounds, identical to Stage~2's protocol. Program summaries are in
\cref{appendix:bkdv-bank}; we highlight three:
\begin{itemize}
  \item \textbf{BKdV-S5} establishes an algebraic identity
        $m_t|_{m=0} = (v-1)(6 v v_x + v_{xxx}) \not\equiv 0$ for sech$^2$
        initial conditions, demonstrating that the m{=}0 (Gardner) reduction
        is \emph{not} a dynamically invariant manifold.
  \item \textbf{BKdV-S6} establishes that the pre-validated stack
        (pseudospectral + 2/3 dealias + RK4) is \emph{quantitatively wrong}
        for bore-like u initial conditions: u develops Gibbs-driven
        non-physical excursions and a usable threshold is
        $\nu_\text{linear} \approx 5\times 10^{-2}$ (or
        $\nu_h \approx 10^{-9}$ for $k^8$ hyperviscosity), \emph{13 orders of
        magnitude} above the BKdV-S4 smooth-soliton ``safe envelope''.
  \item \textbf{BKdV-S7} demonstrates that an initial condition that is
        stable in Gardner-only evolution at amplitude $1.5$ is \emph{not}
        stable in BKdV when initialised at $m_0 = 0$: $v_{\max}$ decays by
        $62.8\%$ over $T{=}10$, with the mode-by-mode amplification of $m$
        predicted by S5's identity to a cosine similarity of $0.94$ against
        the observed spectrum.
\end{itemize}

\paragraph{Dual-pass curation.} For each program a sub-agent runs the three
rounds; two passes of curators then read the complete trace. The
\emph{per-round} pass produces three depth-1 records, capturing each round's
local lesson. The \emph{deep} pass synthesises across all three rounds,
producing one depth-3 record per program containing a
\texttt{synthesised\_diagnosis}, a list of \texttt{ruled\_out\_routes}
(path-level closures), and a \texttt{recommended\_alternative} field that
must not duplicate any ruled-out route. The deep records are the carriers of
multi-round path-closure knowledge.

\paragraph{The expanded bank.} The 30 legacy single-shot entries are retained
(re-tagged as depth-1 records) and supplemented with the 28 new entries (21
depth-1 per-round + 7 depth-3 deep). The final bank \texttt{bank\_v3.A'}
contains 58 entries: 15 positive + 43 negative, with 7 deep-synthesis entries
covering S1--S7. The depth tag is used in \cref{sec:headline-result} to
ablate ``one-round-failure log'' versus ``three-round path-closure''
contributions.

\subsection{Progressive-complexity discipline}
\label{sec:progressive-discipline}

A central methodological problem in this case study is \emph{method lookup}:
when a bank contains validated complex methods (e.g.\ IMEX-CN spectral), a
bank-aware agent at Stage~2 can adopt the complex method at round~1 without
exhibiting research-style reasoning. To prevent this we impose a
non-negotiable discipline on every Stage~2 cell:
\begin{itemize}
  \item \textbf{E1 must be the simplest meaningful baseline}: pseudospectral
        spatial derivatives, no dealiasing, classical explicit RK4, no
        operator splitting, no implicit treatment of stiff terms. Even if the
        bank explicitly endorses a complex stack, the agent must first
        observe baseline failure modes on the actual task.
  \item \textbf{E2 and E3 change at most one major component} over the
        previous experiment. If E2 simultaneously switches the time
        integrator, adds dealiasing, and changes the discretisation, the
        resulting outcome cannot be attributed to any single component. The
        bank serves as guidance for \emph{which} single component to upgrade,
        not as a recipe to bypass baseline.
\end{itemize}
This discipline mirrors normal research practice and is the only way to
measure whether negative knowledge actually narrows an agent's escalation
choices, as opposed to being silently absorbed by an already-stacked method.

\subsection{Stage~2: Four memory conditions on three sub-tasks}
\label{sec:pde-stage2}

\paragraph{Sub-tasks.} \textbf{T-A} soliton stability near the $m=0$
manifold; \textbf{T-B} whether a high-amplitude Gaussian wave packet
decomposes into a soliton train; \textbf{T-C} whether a KdV soliton survives
crossing a Burgers bore. None has a closed-form reference; success is judged
by deterministic phenomenon checks (mass-drift, peak counts, amplitude,
boundedness).

\paragraph{Four conditions.} We extend the original three-condition design
with a fourth condition, \textbf{NegOnly}, which receives only the negative
entries of the bank. The four conditions are: \textbf{NoKB} (no bank),
\textbf{PosOnly} (positive entries only), \textbf{NegOnly} (negative entries
only), \textbf{PosNeg} (full bank). NegOnly tests whether negative knowledge
alone, without positive recommendations, can carry the team to a working
solution---a stringent test of the prescriptive content embedded in negative
entries.

\paragraph{Cells and budget.} The total $3 \times 4 = 12$ cells, each with
$\leq 3$ Research Graph rounds under the progressive-complexity discipline.
Per-task specifications are in \cref{appendix:bkdv-setup}.

\subsection{Physics-aware phenomenon check and one corrected threshold}
\label{sec:eval-correction}

We discovered during analysis that the original T-A phenomenon-check
threshold ($v_{\max}(T)/v_{\max}(0) \geq 0.5$) is \emph{physically
unattainable} for our chosen initial condition. BKdV-S7's quantitative
prediction is that an off-manifold sech$^2$ IC at amplitude $1.5$ decays by
$62.8\%$ over $T{=}10$ as a physical consequence of the m{=}0 manifold not
being dynamically invariant; our T-A IC is more perturbed and decays
similarly. The original threshold mechanically rewarded runs whose
chaotic fragmentation happened to leave a peak above $1.0$ (one out of
14--16 spurious peaks), and punished runs that produced the physically
correct single coherent soliton at the lower amplitude. We correct the
threshold to $\geq 0.25$ and add an explicit single-dominant-peak requirement
(either $n_\text{peaks} = 1$ or top-to-second-peak ratio $> 1.5$). Full
derivation in \cref{appendix:eval-v2}. This correction is itself a finding:
bank-aware research surfaces benchmark mis-calibration that bank-unaware
research cannot, because only bank-aware agents arrive at the physically
correct end-state.

\subsection{Headline result: NegOnly $1/3 \to 3/3$ from one prescriptive negative}
\label{sec:headline-result}

We run Stage~2 first with the 50-entry initial bank (BKdV-S1--S5 only), then
re-run with the expanded 58-entry bank \texttt{bank\_v3.A'} (adding the
8 entries from S6 and S7). Both runs use the physics-aware eval of
\cref{sec:eval-correction}. \Cref{tab:bkdv-headline} reports PASS rates.

\begin{table}[tb]
\centering
\caption{Stage~2 PASS rate (physics-aware eval) under two bank versions.
Bold = changed from initial to expanded.}
\label{tab:bkdv-headline}
\small
\begin{tabular}{l cc cc}
\toprule
& \multicolumn{2}{c}{Initial bank (50)} & \multicolumn{2}{c}{Expanded bank (58)} \\
\cmidrule(lr){2-3}\cmidrule(lr){4-5}
Cond.   & T-A & PASS/3 & T-A & PASS/3 \\
\midrule
NoKB    & $\times$ & 0/3 & $\times$ & 0/3 \\
PosOnly & $\times$ & 1/3 & $\times$ & 1/3 \\
NegOnly & \cmark{} & 1/3 & \cmark{} & \textbf{3/3} \\
PosNeg  & \cmark{} & 3/3 & \cmark{} & 3/3 \\
\bottomrule
\end{tabular}
\end{table}

The largest change is NegOnly: $1/3 \to 3/3$. Adding the BKdV-S6
deep-synthesis entry---which contains a \emph{prescriptive}
\texttt{recommended\_alternative} field (``adopt
$\nu_\text{linear} = 5\times 10^{-2}$ as the default u-side dissipation for
bore-like u-gradient ICs'')---lifts NegOnly to match PosNeg. The agent traces
(\cref{appendix:agent-traces}) show NegOnly cells citing BKdV-S6 by name and
adopting the prescribed coefficient directly; without this single entry,
NegOnly's bank is categorical-vetoes-only and cannot find an actionable
u-side dissipation level within the three-round budget.

\paragraph{Bank content $\gg$ bank size.} The legacy 30-entry single-shot
bank already covered all the same physical phenomena (Gibbs at Burgers
shocks, KdV dispersive stiffness, aliasing under nonlinear products). Adding
20 single-round entries (per-round records from S1--S5) lifted PosNeg only
marginally (already at 3/3) and did not change NoKB. The change comes from
the \emph{type} of content in the 8 newer entries (S6, S7 per-round plus
deep): deep-synthesis records of multi-round path closures with prescriptive
replacement recommendations are categorically more useful than depth-1
single-shot failure logs of the same nominal subjects.

\paragraph{PosNeg robustness.} PosNeg achieves 3/3 in both bank versions. The
full bank lets an agent read prescription (positive) and warning (negative)
together, narrowing the escalation choice from both sides; PosOnly and
NegOnly each see only half the constraint and can be defeated when the other
half is the binding one.

\subsection{Mechanism on T-C: prescription narrows method search}
\label{sec:tc-mechanism}

The bore-soliton sub-task T-C is where the bank's prescriptive content matters
most. Under the initial 50-entry bank, all four conditions fail T-C: NoKB and
PosOnly let the Burgers bore grow Gibbs ringing to $|u|_{\max} > 5$; NegOnly
correctly rejects central FD and Lax--Friedrichs but cannot find a working
u-side dissipation level within three rounds; PosNeg passes (in the original
v2 eval, on a related grid; not under v3.A' progressive-complexity).

Under the expanded 58-entry bank, T-C is solved by three of four conditions
in 2--3 rounds. The decisive entry is BKdV-S6's deep record, which (i)
identifies that the BKdV-S4 ``safe hyperviscosity envelope'' is 13 orders
of magnitude too weak for bore-like u, (ii) prescribes a concrete level
$\nu_\text{linear} = 5\times 10^{-2}$, and (iii) explains why $\varepsilon =
10^{-4}$ linear viscosity (the textbook small-viscosity default) is
underdosed by 2--3 orders of magnitude. The PosNeg cell on T-C at round~2
quotes the BKdV-S6 prescription verbatim in its bank-use rationale while
simultaneously rejecting the BKdV-S4 envelope as ``too weak for bore-like
IC''; full excerpt in \cref{appendix:agent-traces}. The result is a
two-round PASS with $|u|_{\max} \approx 1.3$, an order of magnitude below
the bound.

\subsection{Limitations}
\label{sec:case-study-limitations}

The case study has three honest limitations. First, the initial T-A
phenomenon-check threshold was physically miscalibrated; this was detected
only because the BKdV-S7 stress test produced the correct physical
prediction, but it raises a methodological question of how often benchmark
thresholds for open research questions are similarly mis-calibrated in
unidirectional ways. Second, the positive/negative split is heuristic
(applied to the curator output by failure metadata): we observe one entry
(BKdV-S7 round~2) misclassified as positive due to its
\texttt{failure.recommended\_action = retry} field. The downstream effect is
discussed in the IC-adaptation trap appendix
(\cref{appendix:ic-trap}). Third, Stage~2 sub-agents save only the
\emph{final} solver per cell on disk; per-round candidate code is not
preserved. The research-graph \texttt{research\_state.jsonl} captures the
method description for each Experiment node but is not a byte-exact code
record. Stage~1 sub-agents save full per-round artifacts.

The full bundle---all code, prompts, sub-agent prompts and outputs, both
bank versions, both eval versions, run scripts, the agent trace catalogue,
and per-cell logs---is preserved at \texttt{section4\_reproduce/} for
reproduction.

% =============================================================================
% APPENDIX ADDITIONS (to be appended to existing appendix sections)
% =============================================================================

\section{Stage~1 Stress-Test Program Catalogue}
\label{appendix:bkdv-bank}

\paragraph{Seven multi-round programs.} Each program is a 3-round Research
Graph trace producing 4 NK records (3 per-round single-shot + 1 deep
synthesis): 21 + 7 = 28 new entries to supplement the 30 legacy single-shot
entries from the prior pilot bank. Full prompt templates, agent traces, and
NK record JSONs are in \texttt{section4\_reproduce/stage1\_v3/}.

\begin{itemize}\itemsep0pt
  \item \textbf{BKdV-S1} Numerical-method survey: which combinations stably
        integrate BKdV at amplitudes 1--3 over $T=10$.
  \item \textbf{BKdV-S2} Conservation-law audit: distinguishes
        divergence-form trivialities ($\int u\,dx, \int v\,dx$) from physical
        non-conservations ($\int u v\,dx$, energy candidates).
  \item \textbf{BKdV-S3} IC-family vs.\ coherent structure: smooth-localised
        seeds produce coherent compound structures; broadband seeds at
        amplitude $\geq 0.8$ hit a high-k cascade.
  \item \textbf{BKdV-S4} Numerical-resolution sensitivity: $N_x{=}256$ is
        \emph{sub-converged} for quantitative claims; doubling $N_x$ shifts
        diagnostics by 30--140\%. Defines a hyperviscosity safe envelope for
        smooth ICs ($\nu_h \lesssim 10^{-20}$).
  \item \textbf{BKdV-S5} m{=}0 manifold (non-)invariance: algebraic identity
        and k-selective perturbation response.
  \item \textbf{BKdV-S6} Bore-IC u-viscosity necessity:
        $\nu_\text{linear} \approx 5\times 10^{-2}$ required; BKdV-S4
        envelope underestimates by 13 orders.
  \item \textbf{BKdV-S7} Gardner-stable $\not\Rightarrow$ BKdV-stable:
        $-62.8\%$ $v_{\max}$ decay quantified with cosine-similarity-$0.94$
        mode-by-mode prediction.
\end{itemize}

\begin{table}[tb]
\centering
\caption{Bank \texttt{v3.A'} composition: 58 entries. Legacy = 30 single-shot
records from the pilot bank, retained and re-tagged as depth-1. New = 28
records from S1--S7 (21 per-round depth-1 + 7 deep depth-3).}
\label{tab:bkdv-bank-v3}
\small
\begin{tabular}{lrrr}
\toprule
Source & Positive & Negative & Total \\
\midrule
Legacy single-shot (pilot, depth-1)         & 10 & 20 & 30 \\
S1--S7 per-round (depth-1)                  &  4 & 17 & 21 \\
S1--S7 deep synthesis (depth-3)             &  1 &  6 &  7 \\
\midrule
Total                                       & 15 & 43 & 58 \\
\bottomrule
\end{tabular}
\end{table}

\section{Phenomenon-Check Eval v1 vs v2}
\label{appendix:eval-v2}

\paragraph{T-A threshold derivation.} The original (v1) eval for T-A required
$v_{\max}(T)/v_{\max}(0) \geq 0.5$ and $n_\text{peaks above } 0.5 \geq 1$.
BKdV-S7 quantitative result: a sech$^2$ IC at $A=1.5$ with $u_0 = v_0^2/2$
(so $m_0 = 0$ to machine precision) under BKdV decays from $v_{\max}=1.50$ to
$v_{\max}=0.558$ over $T=10$, a $62.8\%$ amplitude loss. Our T-A IC is
$v = 2\,\text{sech}^2(x+5)$, $u = 0.5\,v^2 + 0.2\,v$, which sits further
off-manifold (initial $m = 0.2 v$, not zero). Bank-unaware runs that landed
in chaotic fragmentation (14--16 peaks each $\sim 1.3$) passed v1; bank-aware
runs that produced the physically correct single coherent peak at
$v_{\max} \approx 0.64$ failed v1. We replace v1 with v2:
\begin{itemize}\itemsep0pt
  \item \textbf{v2 amplitude}: $v_{\max}(T)/v_{\max}(0) \geq 0.25$.
  \item \textbf{v2 coherence}: single-dominant-peak requirement;
        $n_\text{peaks above } 0.4 = 1$ \emph{OR} top-to-second peak ratio
        $> 1.5$. This rejects chaotic fragmentation cases that v1 passed
        mechanically.
\end{itemize}
T-B and T-C eval criteria are unchanged from v1. Source code for both
versions is in \texttt{section4\_reproduce/stage2\_v3/class\_A/eval/}.

\section{IC-Adaptation Trap Experiment}
\label{appendix:ic-trap}

\paragraph{Setup.} The T-B prompt under the expanded bank was modified to
invite IC adaptation from positive bank entries (``you are invited to adapt
this IC if positive knowledge describes a related profile validated to be
stable''). The trap: BKdV-S7's round-1 record (a positive entry) reports a
Gardner-stable IC at $A=1.5$; PosOnly was expected to adopt it as the
T-B initial condition (instead of the reference Gaussian), and fail because
the same IC is BKdV-unstable per BKdV-S7's round-2 record.

\paragraph{Outcome.} No condition adopted the trap IC. Three reasons:
(i)~Agents applied \emph{task-semantic} reasoning---``T-B is Gaussian
decomposition, the bank's sech$^2$ is a different task class''; (ii)~Agents
noted that the T-B reference IC already has $m_0 = -v_0^2/2 \neq 0$, so a
Gardner-stable claim on the m{=}0 manifold does not transfer; (iii)~The
heuristic positive/negative classification leaked BKdV-S7's
round-2 BKdV-breakdown record into the \emph{positive} bank, so PosOnly
in practice saw the warning along with the lure.

\paragraph{What we learn.} The trap result is informative even though the
trap was not triggered: the positive/negative split is more fragile than a
clean binary, and agents read banks holistically rather than treating
``positive entries'' as recommendations to follow without scrutiny. A
cleaner split heuristic (using \texttt{failure.degree} instead of
\texttt{failure.recommended\_action}) would make the trap more discriminating
in future work.

\section{Agent Trace Catalogue}
\label{appendix:agent-traces}

We extract verbatim text from \texttt{research\_state.jsonl} files of
selected cells to illustrate how bank-aware agents use the bank. Each
excerpt is the agent's own field text, written during the run.

\paragraph{Excerpt~1: T-C/PosNeg E2 \texttt{bank\_use\_rationale}.}
The cell passes at round~2 by adopting BKdV-S6's prescriptive $\nu_\text{lin}$
while explicitly rejecting BKdV-S4's hyperviscosity envelope:

\begin{quote}\small
``POSITIVE bank guidance: BKdV-S6 r=3 establishes $\nu_\text{lin} \in
\{10^{-2}, 5\cdot 10^{-2}\}$ both pass; $\nu = 5\cdot 10^{-2}$ yields
$\text{TV}_\text{final}{\sim}9.6$ and $u_\text{min,final} \geq 0$ (best).
NEGATIVE bank guidance: BKdV-S6 deep synthesis explicitly warns that
BKdV-S4's safe envelope ($\nu_h \sim 10^{-22}$ to $10^{-20}$ for $k^8$
hyperviscosity) is 13 orders of magnitude too weak for bore-like IC---we
REJECT transferring that envelope. We also reject the would-be-tempting
alternative of just shrinking dt (E1 already used dt$=10^{-4}$ which is
below post-dealias CFL---dt alone cannot stop the bore cascade per
BKdV-S6 r=2 finding that $\varepsilon=10^{-4}$ viscosity was ineffective).''
\end{quote}

This is the prescription-and-warning pattern in action: the agent reads
\emph{both} which level to use AND which would-be-tempting shortcut to skip.

\paragraph{Excerpt~2: T-A/PosOnly F2 rationale (physics-honest fail).}
Under the expanded bank including BKdV-S7, PosOnly no longer applies the
hyperviscosity ``trick'' that artificially preserves $v_{\max}$ above
threshold; instead it diagnoses the failure as physical:

\begin{quote}\small
``Stack ran cleanly to $T=8$ with mass conserved to machine precision and
$|$fields$|$ bounded. However the soliton is destroyed:
$v_{\max}(T)=0.917 < 1.0$ ($0.46\times$ of initial $2.0$), and $v$ contains
5--6 peaks above $0.7$. This matches the BKdV-S7 r2 physical breakdown of
the m{=}0 manifold under the $-\partial_x(uv)$ coupling---and our IC is more
perturbed ($m_0 = +0.2 v_0 = {\sim}0.4 \text{sech}^2$, not~0). The result
\emph{may} still satisfy the [\ldots]''
\end{quote}

The agent now \emph{cites} the BKdV-S7 mechanism instead of attempting a
numerical workaround. This is a research-character outcome:
bank explanation of why something fails physically discourages numerical
sleight-of-hand.

\paragraph{Excerpt~3: T-C/NegOnly E2 \texttt{bank\_use\_rationale}.}
NegOnly demonstrates that a sufficiently prescriptive negative bank carries
its own actionable content (it does not need positive entries to recommend
a method):

\begin{quote}\small
``Single-component upgrade vs E1: add $\nu \cdot u_{xx}$ on $u$ only
(treated explicitly inside RK4). BKdV-S6 deep synthesis: $\nu=5\cdot 10^{-2}$
empirically passes ($\text{TV}_\text{final} \sim 9.6$,
$u_\text{min,final} \geq 0$) on the EXACT same bore IC used here, while the
alternative hyperviscosity $\nu_h$ must approach the explicit-RK4 stability
ceiling ${\sim}10^{-9}$ (10+ orders above BKdV-S4 `safe envelope' $10^{-22}$).
REJECT BKdV-S4 ladder: $\nu_h \leq 10^{-12}$ is empirically insufficient
($\text{TV}>115$, $u_\text{max}>3$)---would waste an iteration. REJECT
MUSCL/Godunov on $u u_x$: a 2+ component swap from E1 (changes both
discretisation AND adds limiter) violates progressive-complexity [\ldots]''
\end{quote}

This is what we mean by ``prescriptive negative'': the entry both \emph{rules
out a tempting wrong direction} (transferring the BKdV-S4 envelope) and
\emph{names a working level} ($\nu = 5\cdot 10^{-2}$). Without the named
level the bank would still rule out BKdV-S4's $\nu_h \sim 10^{-22}$, but
the agent would have no anchor for what level \emph{does} work.
